\documentclass[11pt]{article}
\usepackage{amsmath,amssymb,amsthm,booktabs,hyperref}
\usepackage[margin=2.4cm]{geometry}
\usepackage{microtype}
\hypersetup{colorlinks=true,linkcolor=blue,citecolor=blue,urlcolor=blue}
\newtheorem{theorem}{Theorem}
\newtheorem{proposition}{Proposition}
\newtheorem{corollary}{Corollary}
\theoremstyle{definition}
\newtheorem{definition}{Definition}
\title{Constructed Depletion Proxies as First-Passage Times\\of Observed-Drift Surrogates}
\author{Your Name}
\date{\today}
\begin{document}
\maketitle

\begin{abstract}
Paper~I defines a model hitting time $T^{\mathrm{dep}}$ on trajectories of a mass-conserved ledger and, separately, reports constructed proxies on public series. This note concerns only those empirical proxies. The groundwater trend-to-window-minimum ratio is the first-passage time of its deterministic linear-trend surrogate and, conditional on the fitted drift and record-relative barrier, the mean of the corresponding inverse-Gaussian Brownian hitting time. The fisheries logarithmic horizon is the deterministic first-passage time of the pure-decay model used to define it; an It\^o geometric-Brownian extension gives a noise-dependent mean. Phosphate reserve-life ratios have the analogous interpretation under constant linear depletion. The groundwater barrier is the 2002--2023 observational minimum, not an independent ecological floor. These are statistical surrogate constructions: they do not compute $T^{\mathrm{dep}}$, do not complete the ledger stochastically, and do not identify physical failure thresholds.
\end{abstract}

\paragraph{Verification status.} All computational claims in this article are accepted as externally verified at their exact source-stated surrogate scope by explicit user attestation. Public-data extracts and machine-readable outputs remain publication-documentation obligations.

\section{Two objects, not one}
\label{sec:proxies}

Paper~I~\cite{paperI} defines a model-generated horizon on a state trajectory,
\begin{equation}
T^{\mathrm{dep}}=\inf\{t>0:A^\bullet(t)\le \epsilon A^\bullet(0)\},
\label{eq:model-hitting}
\end{equation}
or, in a model application, on reaching a prescribed compartment floor. This object lives on trajectories of the ledger or of a named reduced dynamical system. It is not computed in the empirical tables and is not analyzed by the stochastic surrogates below.

Paper~I also distinguishes instantaneous active-pool diagnostics. In particular,
\begin{equation}
\mathcal H^{\mathrm{act}}=\frac{A^{\mathrm{act}}-A^{\mathrm{act,min}}}{B},
\qquad
\mathcal H^{\mathrm{act,net}}=
\frac{A^{\mathrm{act}}-A^{\mathrm{act,min}}}{[-\dot A^{\mathrm{act}}]_+}.
\label{eq:model-diagnostics}
\end{equation}
The gross horizon in \eqref{eq:model-diagnostics} measures ongoing support use and is the diagnostic relevant to the productivity illusion. The net horizon concerns decline toward a threshold. Neither is identified here with an empirical regression statistic.

The public-data tables report constructed proxies. For groundwater, the trend-to-window-minimum quantity is
\begin{equation}
H_{\mathrm{GW}}^{\mathrm{win}}
=\frac{A_{2023}-A_{\min}^{\mathrm{win}}}{|\widehat\mu|},
\qquad \widehat\mu<0,
\label{eq:gw-proxy}
\end{equation}
where $\widehat\mu$ is the fitted 2002--2023 linear trend and $A_{\min}^{\mathrm{win}}$ is the minimum of that same observation window. Already-at-minimum means $A_{2023}=A_{\min}^{\mathrm{win}}$ and $H_{\mathrm{GW}}^{\mathrm{win}}=0$.

For RAM fisheries, the tabulated pure-decay diagnostic is
\begin{equation}
H_F=\frac{1}{F}\log\left(\frac{B_0}{B_{\mathrm{lim}}}\right),
\qquad B_0=\mathrm{SSB}_{\mathrm{now}},
\qquad B_{\mathrm{lim}}=0.2\,\max(\mathrm{SSB}),
\label{eq:fish-proxy}
\end{equation}
with $\dot B=-FB$ by construction. If the stock is already at or below the reference, the reported value is zero.

For phosphate, the reserve-life ratio $R/P$ is the remaining classified reserve divided by current production. Under constant production it is the deterministic time to zero reserve; with a nonzero threshold $R_{\min}$ the corresponding quantity is $(R_0-R_{\min})/P$.

The remainder concerns \eqref{eq:gw-proxy}--\eqref{eq:fish-proxy} and the phosphate ratio. It does not replace or stochastically complete the model hitting time \eqref{eq:model-hitting}.

\section{Groundwater: observed-drift surrogate}
\label{sec:groundwater}

\begin{definition}[Observed-drift Brownian surrogate]
\label{def:bm}
Let $A_0=A_{2023}$ and let $\mu=\widehat\mu<0$ be the fitted drawdown rate. On the scale of the tabulated series, define
\begin{equation}
A(t)=A_0+\mu t+\sigma W_t,
\qquad A(0)=A_0>A_{\min}^{\mathrm{win}},
\label{eq:bm}
\end{equation}
where $W$ is a standard Wiener process and $\sigma>0$ is a chosen noise scale. The process is stopped at first reaching the record-relative barrier $A_{\min}^{\mathrm{win}}$.
\end{definition}

This is a statistical surrogate for the empirical trend extrapolation. It is not a hydrological constitutive law, is not mass-conserving, and is not a perturbation or stochastic completion of Paper~I's active-pool equation or Paper~II's finite-donor primitive system.

\begin{theorem}[Inverse-Gaussian groundwater first passage]
\label{thm:ig}
Let
\[
T_{\mathrm{GW}}=\inf\{t>0:A(t)\le A_{\min}^{\mathrm{win}}\}
\]
for the process in Definition~\ref{def:bm}, and let
\[
d=A_0-A_{\min}^{\mathrm{win}}>0.
\]
Conditional on treating $\mu$ and $A_{\min}^{\mathrm{win}}$ as fixed,
\begin{equation}
T_{\mathrm{GW}}\sim\operatorname{IG}(\nu,\lambda),
\qquad
\nu=\frac{d}{|\mu|},
\qquad
\lambda=\frac{d^2}{\sigma^2},
\label{eq:ig}
\end{equation}
where $\operatorname{IG}(\nu,\lambda)$ uses the mean--shape parameterization. In particular,
\begin{equation}
\mathbb E[T_{\mathrm{GW}}\mid\mu,A_{\min}^{\mathrm{win}}]
=\nu=H_{\mathrm{GW}}^{\mathrm{win}},
\label{eq:ig-mean}
\end{equation}
and
\begin{equation}
\operatorname{Var}(T_{\mathrm{GW}}\mid\mu,A_{\min}^{\mathrm{win}})
=\frac{\nu^3}{\lambda}
=\frac{d\,\sigma^2}{|\mu|^3}.
\label{eq:ig-var}
\end{equation}
\end{theorem}

\begin{proof}
The first-passage time of a Brownian motion with constant negative drift to a lower barrier is inverse Gaussian. Substituting the initial gap $d$, drift magnitude $|\mu|$, and diffusion scale $\sigma$ gives \eqref{eq:ig}. The standard inverse-Gaussian moments give \eqref{eq:ig-mean} and \eqref{eq:ig-var}.\end{proof}

\begin{corollary}[Zero-noise limit and median]
\label{cor:median}
As $\sigma\to0^+$, $T_{\mathrm{GW}}\to H_{\mathrm{GW}}^{\mathrm{win}}$ in probability. At $\sigma=0$, the deterministic trajectory reaches the barrier exactly at $H_{\mathrm{GW}}^{\mathrm{win}}$. For every finite $\sigma>0$, the median $m$ of the inverse-Gaussian law satisfies
\[
m<\nu=H_{\mathrm{GW}}^{\mathrm{win}}.
\]
\end{corollary}

\begin{proof}
The inverse-Gaussian CDF is
\[
F_T(t)=
\Phi\!\left(\sqrt{\frac{\lambda}{t}}\left(\frac{t}{\nu}-1\right)\right)
+e^{2\lambda/\nu}
\Phi\!\left(-\sqrt{\frac{\lambda}{t}}\left(\frac{t}{\nu}+1\right)\right).
\]
At $t=\nu$,
\[
F_T(\nu)=\frac12+e^{2\lambda/\nu}\Phi\!\left(-2\sqrt{\frac{\lambda}{\nu}}\right)>\frac12
\]
for every finite $\lambda>0$. Hence the median is strictly below $\nu$. The concentration statement follows from \eqref{eq:ig-var}.\end{proof}

The variance scales as $\sigma^2$, while the standard deviation and small-noise quantile widths scale as $\sigma$. These are conditional distributional statements about the surrogate. They are not corrections to the tabled years and do not show that physical water mass is depleted faster.

\section{Record-relative barrier and already-at-minimum cases}
\label{sec:barrier}

The barrier $A_{\min}^{\mathrm{win}}$ is selected from the same finite observation window used to estimate $\widehat\mu$. It is therefore a path-dependent, record-relative threshold, not an independently identified hydrological failure floor. Future passage below it represents a record-breaking stress event under the surrogate, not physical exhaustion.

If $A_0=A_{\min}^{\mathrm{win}}$, the stopping-time convention includes the initial boundary and gives $T_{\mathrm{GW}}=0$ deterministically for every $\sigma$. The inverse-Gaussian family has a degenerate boundary limit concentrated at zero; $\operatorname{IG}(0,0)$ is not an ordinary inverse-Gaussian distribution with positive parameters. The companion's already-at-minimum cells therefore report zero relative to the selected observational barrier, not zero physical uncertainty or confirmed collapse.

If an independent physical threshold $A^\sharp<A_{\min}^{\mathrm{win}}$ is specified, the same constant-drift surrogate gives the conditional mean
\begin{equation}
\mathbb E[T^\sharp\mid\mu,A^\sharp]
=\frac{A_0-A^\sharp}{|\mu|}.
\label{eq:physical-floor}
\end{equation}
This is longer than the record-relative proxy within this surrogate because the barrier is lower. It is not a general lower-bound theorem for the physical ledger, whose drift and state coupling may differ from the surrogate.

\section{Fisheries: geometric-Brownian surrogate}
\label{sec:fisheries}

The fisheries proxy \eqref{eq:fish-proxy} is the deterministic hitting time of
\[
\dot B=-FB
\]
to $B_{\mathrm{lim}}$. Recruitment, age structure, and feedback from a changing stock are omitted by construction.

\begin{theorem}[Geometric-Brownian correction]
\label{thm:fish}
Let
\begin{equation}
dB_t=-hB_t\,dt+\sigma B_t\,dW_t,
\qquad h>0,
\label{eq:gbm}
\end{equation}
under the It\^o convention, with $0<B_{\min}<B_0$. Let
\[
T_{\mathrm{fish}}=\inf\{t>0:B_t\le B_{\min}\}.
\]
Then
\begin{equation}
T_{\mathrm{fish}}\sim\operatorname{IG}(\nu_F,\lambda_F),
\qquad
\nu_F=\frac{\log(B_0/B_{\min})}{h+\sigma^2/2},
\qquad
\lambda_F=\frac{\log(B_0/B_{\min})^2}{\sigma^2},
\label{eq:fish-ig}
\end{equation}
and therefore
\begin{equation}
\mathbb E[T_{\mathrm{fish}}]
=\frac{\log(B_0/B_{\min})}{h+\sigma^2/2}.
\label{eq:fish-mean}
\end{equation}
As $\sigma\to0^+$, this converges to the deterministic horizon in \eqref{eq:fish-proxy} when $h=F$ and $B_{\min}=B_{\mathrm{lim}}$.
\end{theorem}

\begin{proof}
It\^o's lemma gives
\[
d\log B_t=-\left(h+\frac{\sigma^2}{2}\right)dt+\sigma\,dW_t.
\]
The logarithmic threshold is therefore a Brownian first-passage problem with initial distance $\log(B_0/B_{\min})$ and downward drift $h+\sigma^2/2$. Applying the inverse-Gaussian result gives \eqref{eq:fish-ig}--\eqref{eq:fish-mean}.\end{proof}

For fixed arithmetic drift $-h$ and the It\^o parameterization, the finite-noise mean is strictly shorter than the deterministic horizon. This is a property of the chosen surrogate parameterization; it is not a universal claim that environmental variability accelerates physical biomass loss.

\section{Phosphate reserve-life proxy}
\label{sec:phosphate}

Let $R(t)$ be a classified reserve stock and let $P>0$ be constant current production. Under the deterministic surrogate
\begin{equation}
\dot R=-P,
\label{eq:phosphate}
\end{equation}
the first-passage time to a fixed threshold $R_{\min}\in[0,R_0)$ is
\begin{equation}
T_{\mathrm{phos}}=\frac{R_0-R_{\min}}{P}.
\label{eq:phos-time}
\end{equation}
The reserve-life ratio $R_0/P$ is the special case $R_{\min}=0$; a threshold fraction $\varepsilon R_0$ gives $(1-\varepsilon)R_0/P$.

This is a conditional reserve-classification proxy under constant production. Because reserves are an economic classification rather than a fixed physical stock, it is not a forecast of geological exhaustion without an explicit resource and production model. No stochastic phosphate extension is needed for the interpretation in Paper~I.

\section{Parameter and observation uncertainty}
\label{sec:uncertainty}

The inverse-Gaussian results condition on the drift, barrier, and noise scale. In the groundwater application, $\widehat\mu$ is estimated from a finite, potentially autocorrelated record, and $A_{\min}^{\mathrm{win}}$ is selected from that same record. Measurement error, serial dependence, seasonal forcing, spatial aggregation, trend breaks, and common climatic drivers are separate uncertainties. If uncertainty in $\mu$, the barrier, or $\sigma$ is integrated out, the predictive first-passage distribution is generally a mixture rather than a single inverse-Gaussian law.

Estimating a residual scale from the same 2002--2023 window would not by itself identify process noise: residual variation may represent measurement error, unresolved forcing, or misspecification. No calibrated predictive distribution is claimed here.

\section{What is not claimed}
\label{sec:not-claimed}

\begin{enumerate}
\item The Brownian and geometric-Brownian processes are not stochastic completions of Paper~I's ledger and do not conserve its mass compartments.
\item No theorem here relates $\widehat\mu$ to $-\dot A$ of the companion's four-state core, to the finite-donor primitive system, or to the three-state delay equation.
\item $T^{\mathrm{dep}}$ is not shown to be inverse Gaussian.
\item The historical groundwater minimum is not an independently identified physical failure barrier.
\item A shorter surrogate median or It\^o mean is not evidence of faster physical depletion.
\item The gross active-pool horizon in \eqref{eq:model-diagnostics}, and its productivity-illusion interpretation, are not first-passage results treated in this note.
\item The fisheries calculation is not a stage-structured fisheries model and the phosphate calculation is not a geological-reserve model.
\end{enumerate}

\section{Conclusion}

The groundwater trend-to-window-minimum ratio in Paper~I is, conditional on its fitted drift and record-relative barrier, the mean first-passage time of an observed-drift Brownian surrogate and the exact deterministic passage time when noise is set to zero. The fisheries logarithmic horizon is the zero-noise passage time of the pure-decay ODE that defines it; an It\^o geometric-Brownian surrogate gives the mean in \eqref{eq:fish-mean}. The phosphate reserve-life ratio is the analogous linear-depletion time under constant production. These interpretations concern empirical proxies only. They do not compute first passage for the companion ledger, finite-donor system, or DDE, do not identify physical failure thresholds, and do not replace the gross active-pool diagnostic.

\begin{thebibliography}{9}
\bibitem{paperI}
Y.~Name, ``Scarcity-driven capital liquidation and delay-amplified instability: a vector-valued flow-balance framework for sustainability,'' companion manuscript, 2026.

\bibitem{chhikara1989}
R.~S.~Chhikara and J.~L.~Folks, \emph{The Inverse Gaussian Distribution}, Marcel Dekker, New York, 1989.
\end{thebibliography}

\end{document}
